\documentclass[11pt]{amsart}
\usepackage[T1]{fontenc}
\usepackage{lmodern,microtype,amsmath,amssymb,amsthm,mathtools}
\usepackage[margin=1.1in]{geometry}
\usepackage[hidelinks]{hyperref}
\usepackage{xurl}

\newtheorem{theorem}{Theorem}[section]
\newtheorem{lemma}[theorem]{Lemma}
\newtheorem{corollary}[theorem]{Corollary}
\theoremstyle{definition}
\newtheorem{definition}[theorem]{Definition}
\theoremstyle{remark}
\newtheorem{remark}[theorem]{Remark}

\newcommand{\N}{\mathbb N}
\newcommand{\C}{\mathbb C}
\DeclareMathOperator{\Sym}{Sym}
\DeclareMathOperator{\Stab}{Stab}

\hypersetup{pdftitle={A homogeneous space with no nonconstant quasiharmonic functions},pdfauthor={Yitzchak Shmalo}}

\title[No nonconstant quasiharmonic functions on a homogeneous space]{A homogeneous space with no nonconstant\\ quasiharmonic functions}
\author{Yitzchak Shmalo}
\address{Einstein Institute of Mathematics, The Hebrew University of Jerusalem, Jerusalem, Israel}

\subjclass[2020]{60J50, 60B15, 20B07, 31C05}
\keywords{random walk on groups, harmonic functions, quasiharmonic functions, homogeneous space, finitary symmetric group}

\begin{document}

\begin{abstract}
Bj\"orklund asked for a countable measured group $(G,\mu)$ and a subgroup $H$ of infinite index such that the homogeneous space $G/H$ carries no nonconstant left $\mu$-harmonic or $\mu$-quasiharmonic functions, bounded or not. We give such an example. The group is the group of finitely supported permutations of $\N_0=\{0,1,2,\dots\}$, the subgroup is the stabiliser of $0$, and $\mu$ is any probability measure that charges every transposition $(0\ n)$, $n\ge1$, and nothing else. The induced walk on $G/H\cong\N_0$ is a star with infinitely many leaves. The equation at a leaf $k$ forces a quasiharmonic function with defect $c$ to satisfy $u(k)=u(0)-c/\mu((0\ k))$, and the equation at the centre then forces $c=0$ and $u$ constant. The argument uses only the convergence of the single series that defines the equation at the centre.
\end{abstract}

\maketitle

\section{Introduction}

A \emph{measured group} is a pair $(G,\mu)$ where $G$ is a countable group and $\mu$ is a probability measure on $G$ whose support generates $G$ as a semigroup. Write $\check\mu(s)=\mu(s^{-1})$. Following~\cite[Sections~1 and~4.1]{Bj}, a complex-valued function $\phi$ on $G$ is \emph{left $\mu$-harmonic} if
\begin{equation}\label{eq:harm}
\int_G \phi(sg)\,d\check\mu(s)=\phi(g)\qquad(g\in G),
\end{equation}
and \emph{left $\mu$-quasiharmonic} if there is a constant $\ell(\phi)\in\C$ with
\begin{equation}\label{eq:qharm}
\int_G \phi(sg)\,d\check\mu(s)=\phi(g)+\ell(\phi)\qquad(g\in G),
\end{equation}
in both cases for $\mu$-integrable $\phi$. Left translation commutes with right translation, so for a subgroup $H<G$ these equations make sense for functions on the space $G/H$ of left cosets, that is, for right $H$-invariant functions on $G$. Constant functions are harmonic, and harmonic functions are quasiharmonic with $\ell=0$.

Bj\"orklund studies harmonic and quasiharmonic functions on homogeneous spaces in~\cite[Appendix~II]{Bj}. His Problem~2 in~\cite[Appendix~II, Section~7.1]{Bj} reads: construct a countable measured group $(G,\mu)$ and an infinite index subgroup $H<G$ such that $G/H$ admits no non-constant left (quasi-)$\mu$-harmonic functions whatsoever. The text before it remarks that for general quotient spaces there could well be simple counterexamples. The following theorem gives one.

\begin{theorem}\label{thm:main}
Let $G=\Sym_{\mathrm{fin}}(\N_0)$ be the group of permutations of $\N_0$ that move only finitely many points, let $H=\Stab_G(0)$, and let $\mu$ be a probability measure on $G$ with $\mu((0\ n))=p_n>0$ for every $n\ge1$ and $\sum_{n\ge1}p_n=1$. Then $(G,\mu)$ is a measured group, $\mu$ is symmetric, $H$ has infinite index in $G$, and every left $\mu$-quasiharmonic function on $G/H$ is constant. In particular every left $\mu$-harmonic function on $G/H$ is constant.
\end{theorem}

A concrete choice is $p_n=2^{-n}$. The proof in Section~\ref{sec:proof} shows a little more: a function $u$ on $G/H$ satisfying the quasiharmonic equation is constant as soon as the one series defining its left side at the coset $H$ converges; absolute convergence is not needed (Theorem~\ref{thm:star}). The group $G$ is not finitely generated and $\mu$ has infinite support. Both features are used: the centre of the star has infinitely many neighbours, and this is what rules out a nonzero defect. The problem with $G$ finitely generated or $\mu$ finitely supported is not treated here.

\section{The measured group and its quotient}

For $n\ge1$ write $\tau_n=(0\ n)$, and let $\mu$ be as in Theorem~\ref{thm:main}.

\begin{lemma}\label{lem:group}
The group $G$ is countable, $\check\mu=\mu$, the support of $\mu$ generates $G$ as a semigroup, and the map $gH\mapsto g(0)$ is a bijection from $G/H$ onto $\N_0$ satisfying $(sg)(0)=s(g(0))$ for $s,g\in G$. In particular $[G:H]=\infty$.
\end{lemma}

\begin{proof}
Every element of $G$ lies in $\Sym(\{0,1,\dots,m\})$ for some $m$, so $G$ is a countable union of finite groups. Each $\tau_n$ is an involution, so $\check\mu=\mu$. For distinct $i,j\ge1$ one has $\tau_i\tau_j\tau_i=(i\ j)$, since conjugating $(0\ j)$ by $(0\ i)$ replaces $0$ by $i$. Hence the semigroup generated by the $\tau_n$ contains every transposition of $\N_0$, and every element of $G$ is a finite product of transpositions. The map $gH\mapsto g(0)$ is well defined because $H$ fixes $0$; it is injective because $g(0)=g'(0)$ means $g^{-1}g'\in H$; and it is onto because $\tau_n(0)=n$. The identity $(sg)(0)=s(g(0))$ is the definition of composition.
\end{proof}

Identify $G/H$ with $\N_0$ through Lemma~\ref{lem:group}, and let $u\colon\N_0\to\C$ correspond to the right $H$-invariant function $\tilde u(g)=u(g(0))$ on $G$. Since $\check\mu=\mu$ and $\tau_n^{-1}=\tau_n$, the left side of~\eqref{eq:qharm} for $\tilde u$ at $g$ equals $(Pu)(g(0))$, where
\[
(Pu)(x)=\sum_{n\ge1}p_n\,u(\tau_n x)\qquad(x\in\N_0).
\]
Because $\tau_n0=n$, while for $k\ge1$ one has $\tau_nk=0$ if $n=k$ and $\tau_nk=k$ otherwise,
\begin{align}
(Pu)(0)&=\sum_{n\ge1}p_n\,u(n),\label{eq:centre}\\
(Pu)(k)&=p_k\,u(0)+(1-p_k)\,u(k)\qquad(k\ge1).\label{eq:leaf}
\end{align}
This is the transition operator of the walk on the infinite star with centre $0$ that jumps from the centre to the leaf $n$ with probability $p_n$ and from the leaf $k$ to the centre with probability $p_k$, staying put otherwise. The series in~\eqref{eq:leaf} has only two distinct values of $u$ in it and converges absolutely for every $u$; the series~\eqref{eq:centre} is the only one whose convergence depends on $u$.

\section{Proof of the theorem}\label{sec:proof}

\begin{theorem}\label{thm:star}
Let $u\colon\N_0\to\C$ and $c\in\C$. Suppose that the series $\sum_{n\ge1}p_nu(n)$ converges, in the order $n=1,2,\dots$, and that $Pu=u+c$ on $\N_0$. Then $c=0$ and $u$ is constant.
\end{theorem}

\begin{proof}
For $k\ge1$, equation~\eqref{eq:leaf} and $(Pu)(k)=u(k)+c$ give $p_k\,u(0)+(1-p_k)\,u(k)=u(k)+c$, that is,
\begin{equation}\label{eq:forced}
p_k\bigl(u(0)-u(k)\bigr)=c,\qquad u(k)=u(0)-\frac{c}{p_k}.
\end{equation}
Hence the $n$-th term of the series in~\eqref{eq:centre} is $p_nu(n)=p_nu(0)-c$. Since $\sum_np_n=1$, we have $p_n\to0$, so the terms tend to $-c$. The terms of a convergent series tend to $0$, so $c=0$, and then~\eqref{eq:forced} gives $u(k)=u(0)$ for every $k\ge1$.
\end{proof}

\begin{proof}[Proof of Theorem~\ref{thm:main}]
Lemma~\ref{lem:group} gives the statements about $G$, $\mu$ and $H$. If $\phi$ is a left $\mu$-quasiharmonic function on $G/H$, then $u=\phi$ viewed on $\N_0$ satisfies $Pu=u+\ell(\phi)$, and the left side of~\eqref{eq:qharm} at $g=e$ is the convergent series $\sum_np_nu(n)$. Theorem~\ref{thm:star} shows that $\ell(\phi)=0$ and $u$ is constant. Harmonic functions are the case $\ell=0$.
\end{proof}

\begin{remark}
The value of the defect is not normalised anywhere: at the coset $H$ the equation reads $\ell(\phi)=\sum_np_nu(n)-u(0)$, and nonzero constants are harmonic. The conclusion is that nonconstant solutions do not exist, not that every solution vanishes.
\end{remark}

\begin{remark}
The same proof applies to $\Sym_{\mathrm{fin}}(X)$ and the stabiliser of a point $x_0$, for any countably infinite set $X$, with $\mu$ charging every transposition $(x_0\ x)$, $x\ne x_0$. Positivity of every $p_n$ is used: if $p_k=0$, then $\tau_k$ is not in the support and $u(k)$ may be changed freely, so the support no longer generates $G$ and nonconstant harmonic functions appear.
\end{remark}

\section*{Computations}
The file \texttt{code/checks.py} in~\cite{code} (Python~3, standard library, exact integer and rational arithmetic; its output is \texttt{code/checks\_output.txt}) checks on finite truncations the identities used above: $\tau_i\tau_j\tau_i=(i\ j)$ on $\{0,\dots,N\}$ for $N\le8$; that positive words in $\tau_1,\dots,\tau_N$ give all $(N+1)!$ permutations of $\{0,\dots,N\}$ for $N\le6$; that $gH\mapsto g(0)$ is well defined, injective on cosets and equivariant for $N\le5$; the rows~\eqref{eq:centre} and~\eqref{eq:leaf} on $\{0,\dots,12\}$; and the identities~\eqref{eq:forced} and $p_ku(k)=p_ku(0)-c$ for $p_k=2^{-k}$, $k\le60$. The proofs above do not depend on this file.

\section*{Acknowledgements and funding}
The research presented in this paper was supported by the European Research
Council (ERC) under the European Union's Horizon Europe research and
innovation programme (grant agreement No.~101041711), by the Simons
Foundation as part of the Collaboration on the Mathematical and Scientific
Foundations of Deep Learning, by Heights Labs, by Convex Nexus Capital, by
the Israel Science Foundation (grant number 2258/19), and by the Israel
Science Foundation (ISF Grant 4101/25). Computational resources were
provided in part by Google Cloud through its education credits program
(credit 530550103).

\section*{Declaration of competing interest}
The author holds equity in Heights Labs and in Convex Nexus Capital and
has a management role at Convex Nexus Capital. These entities are included
in the funding declaration above.

\section*{Declaration of generative AI and AI-assisted technologies}
The paper comes out of Math Brain, an ongoing project on the use of AI systems for research mathematics. In September 2026 an OpenAI GPT model, running as a research task of that project, proposed the construction and the argument. In October 2026 Claude Opus~5.5 (Anthropic), working in Claude Code, checked the proof against the source, wrote the computer checks of the repository~\cite{code}, and wrote the paper. The author directed the work and is responsible for the content of the paper and for decisions concerning its dissemination.

\section*{Data availability}
No data were used. The code is available at\\ \url{https://github.com/yspennstate/bjorklund-star-quotient}.

\begin{thebibliography}{9}
\bibitem{Bj} M.~Bj\"orklund, \emph{Five remarks about random walks on groups}, preprint, 2014, arXiv:1406.0763v2.
\bibitem{code} Y.~Shmalo, \emph{Code for this paper}, \url{https://github.com/yspennstate/bjorklund-star-quotient}.
\end{thebibliography}

\end{document}
